\section{Анализ предметной области}
\subsection{Понятие и принципы онлайн аренды вещей}

Электронная коммерция является частью цифровой экономики и охватывает все финансовые и торговые операции, совершаемые с использованием компьютерных сетей, а также связанные с ними бизнес-процессы. Онлайн-аренда вещей представляет собой отдельное направление электронной коммерции, в рамках которого через интернет-платформы на платной основе временно передаётся право пользования материальными объектами.

Виды электронной коммерции:
\begin{itemize}
\item P2P (peer-to-peer) – между пользователями;
\item B2B (business-to-business) – между предприятиями;
\item B2C (business-to-consumer) – между предприятием и потребителем;
\item C2C (consumer-to-consumer) – между потребителями;
\item B2G (business-to-government) – между предприятием и государством;
\end{itemize}

Онлайн-бронирование — это бронирование через интернет в интерактивном режиме. Термин используется применительно к бронированию гостиничных номеров, билетов, мест в ресторанах и театрах и т.д.

Данный формат взаимодействия даёт существенные преимущества как арендодателям, так и арендаторам:
\begin{enumerate}
\item Экономическая выгода: арендаторы экономят на покупке редко используемых или дорогих товаров, а владельцы получают доход от простаивающего имущества.
\item Удобство и доступность: поиск, бронирование и оплата возможны круглосуточно из любого места с выходом в интернет, что сокращает время на поиск и сравнение предложений.
\item Прозрачность взаимодействия: платформы позволяют формировать рейтинги, оставлять отзывы и проверять репутацию участников, что повышает доверие при сделках.
\item Удобство и доступность: поиск, бронирование и оплата осуществляются круглосуточно из любой точки с доступом к интернету, что экономит время на поиск и сравнение предложений.
\end{enumerate}

Вместе с тем существуют характерные риски и ограничения:
\begin{enumerate}
\item Контроль качества и соответствия: отсутствует возможность физического осмотра вещи до бронирования, возможны расхождения между описанием, фотографиями и реальным состоянием.
\item Логистика и сохранность: возникают дополнительные расходы и риски при передаче, повреждении или утрате имущества во время доставки или использования.
\item Логистика и сохранность: возникают дополнительные расходы и риски при передаче, повреждении или утрате имущества в процессе доставки или эксплуатации.
\item Юридические и финансовые споры: затруднено возмещение ущерба, возврат залогов и разрешение конфликтов при отсутствии встроенных механизмов страхования.
\item Кибербезопасность: платформа подвержена рискам утечки персональных данных, мошеннических действий и нарушения прав пользователей.
\end{enumerate}

\subsection{Классификация онлайн площадок для аренды}
Рынок онлайн-площадок для аренды вещей неоднороден и включает сервисы, которые можно классифицировать по нескольким ключевым признакам: типу предлагаемых товаров, модели взаимодействия с пользователями, а также типу арендодателей. Понимание этой классификации необходимо для точного позиционирования проектируемой площадки и определения её целевой аудитории.
\subsubsection{По типу предлагаемых товаров}
По специализации онлайн-площадки для аренды делятся на универсальные и нишевые. Универсальные площадки предлагают аренду широкого спектра товаров: от инструментов и электроники до одежды и спортивного инвентаря. Их главное преимущество — широкая аудитория и большое количество предложений, однако им сложнее создавать специализированные инструменты для каждой отдельной категории. Узкоспециализированные площадки, напротив, сосредоточены на одной или нескольких смежных категориях товаров. Такие сервисы выигрывают за счёт глубокого понимания потребностей своей аудитории, предоставления специализированного функционала и формирования более сплочённого сообщества пользователей с высоким уровнем экспертизы в конкретной области.
\subsubsection{По модели взаимодействия с пользователями}
По способу монетизации и степени вовлечённости в транзакцию выделяют три модели онлайн-площадок для аренды. Первая — модель классифайда, когда платформа выступает только доской объявлений, а арендодатель и арендатор самостоятельно договариваются об оплате и условиях сделки. Вторая — модель с полным сопровождением сделки, при которой платформа берёт на себя обработку платежей, страхование и урегулирование споров, взимая комиссию с каждой транзакции. Именно на эту модель ориентируется разрабатываемая система. Третья — подписочная модель, предполагающая ежемесячную плату за неограниченный доступ к аренде вещей из определённого набора, что особенно востребовано в сегменте детских товаров и одежды.
\subsubsection{По типу арендодателей}
По типу участников сделки выделяют две основные модели онлайн-площадок для аренды. Первая — C2C (Consumer-to-Consumer), где арендодателями являются обычные частные лица, сдающие свои личные вещи во временное пользование другим людям. Вторая — B2C (Business-to-Consumer), где в роли арендодателей выступают профессиональные компании проката, использующие платформу как дополнительный канал сбыта своих услуг.

\subsection{Основные бизнес-процессы системы аренды}
\subsubsection{Регистрация и верификация пользователей}
Для минимизации рисков, связанных с мошенничеством или порчей имущества, платформа должна иметь многоуровневую систему регистрации. Пользователь может выступать в двух ролях: арендодатель и арендатор. Для получения доступа к сервису пользователь проходит обязательную верификацию: подтверждение email, а в некоторых случаях — предоставление данных банковской карты и подтверждение личности через государственные сервисы (например, Госуслуги).
\subsubsection{Публикация и поиск предложений}
Арендодатель создаёт карточку товара (объявление), в которой указывает название вещи, её категорию, подробное описание, фотографии, стоимость аренды и доступные даты. Платформа, в свою очередь, предоставляет систему категорий, фильтров и поиска.
\subsubsection{Бронирование и совершение сделки}
Ключевой этап, который платформа автоматизирует с помощью встроенного календаря бронирования. Это позволяет избежать задвоения заказов. Арендатор выбирает период аренды, и система автоматически рассчитывает итоговую стоимость. Оплата всегда проходит через платформу, которая временно замораживает средства до подтверждения обеими сторонами. Это гарантирует, что арендодатель получит деньги, а арендатор — вещь в надлежащем состоянии.
\subsubsection{Бронирование и совершение сделки}
В C2C-сервисах важнейшим компонентом является система рейтингов и отзывов. После завершения сделки арендатор оценивает состояние вещи и оперативность владельца. Рейтинг напрямую влияет на видимость объявлений в поиске и готовность пользователей заключать сделки.
\subsection{Компоненты программно-информационной системы}
Программно-информационная система онлайн-площадки аренды представляет собой комплекс взаимодействующих программных модулей, обеспечивающих хранение, обработку и предоставление информации всем участникам платформы. Рассмотрим основные компоненты.
\subsubsection{Модуль управления пользователями}
Обеспечивает регистрацию, аутентификацию и авторизацию пользователей. Хранит профили с историей транзакций, рейтингами, отзывами и документами верификации.
\subsubsection{Модуль каталога и поиска}
Обеспечивает создание, редактирование и публикацию объявлений. Реализует полнотекстовый поиск с фильтрацией по категориям и дате доступности.
\subsubsection{Модуль бронирования и календаря}
Управляет доступностью товаров во времени, предотвращает двойное бронирование. Ведёт календарь занятости для каждого объявления. Отправляет уведомления арендодателю о поступивших запросах на бронирование и арендатору — о подтверждении или отклонении.
\subsubsection{Модуль генерации договоров}
Встроенный инструмент для создания договоров аренды между участниками сделки на основе их данных из учётных записей на площадке и данных из карточек объявлений. Данный модуль реализует формирование, а также последующее подписание и расторжение договора аренды двумя сторонами с использованием средств ЭЦП на определённый срок.
\subsubsection{Модуль коммуникации}
Встроенный мессенджер обеспечивает безопасное общение между арендатором и арендодателем без раскрытия личных контактных данных. Система push-уведомлений и email-рассылок информирует пользователей о ключевых событиях: бронировании, оплате, возврате, заявках на аренду и прочих уведомлениях в чатах.
\subsubsection{Модуль аналитики и отчётности}
Предоставляет арендодателям дашборд с данными о доходах, популярности объявлений. Доступна сводная статистика: оборот, количество сделок. Арендаторам предоставляется дашборд с количеством объявлений и историей их аренды.

\subsection{Правовое регулирование сферы онлайн-аренды}

Деятельность онлайн-площадки по сдаче вещей в аренду регулируется рядом нормативно-правовых актов Российской Федерации. Основой является Гражданский кодекс РФ, а именно глава 34 «Аренда», которая определяет общие положения о договоре аренды, права и обязанности арендодателя и арендатора, ответственность за недостатки сданного в аренду имущества.\cite{zakon2}

Поскольку арендатор использует вещь для личных, семейных или домашних нужд, на отношения между участниками сделки распространяются положения Закона РФ «О защите прав потребителей» от 07.02.1992 № 2300-1. Согласно данному закону, арендатор имеет право на получение полной и достоверной информации об услуге, на отказ от исполнения договора в любое время при условии оплаты арендодателю фактически понесённых расходов, а также на возмещение убытков, причинённых вследствие недостатков оказанной услуги. Арендодатель, в свою очередь, обязан предоставить вещь в состоянии, соответствующем условиям договора аренды и назначению имущества, и несёт ответственность за недостатки, обнаруженные в течение срока аренды.\cite{zakon1}

Платформа, являясь оператором персональных данных, обязана соблюдать требования Федерального закона № 152-ФЗ «О персональных данных» от 27.07.2006 г. Вся собираемая информация о пользователях (паспортные данные, контактная информация, платёжные реквизиты) должна обрабатываться на законном основании, а также надёжно храниться и защищаться от несанкционированного доступа.\cite{zakon4} Для реализации функций безопасной сделки платформа интегрируется с платёжными сервисами, сертифицированными по стандарту PCI DSS или ISO-20022, что обеспечивает безопасность финансовых транзакций.

Важным аспектом правового регулирования является налогообложение доходов арендодателей. Арендодатели — физические лица, не зарегистрированные как индивидуальные предприниматели, могут применять специальный налоговый режим «Налог на профессиональный доход» в соответствии с Федеральным законом № 422-ФЗ от 27.11.2018 г.\cite{zakon3}

\subsection{Анализ существующих решений и аналогов}
Для улучшения качества собственного продукта следует провести анализ уже существующих решений. Сфера онлайн-бронирования с каждым годом развивается и расширяется. Наиболее востребованными в данном сегменте являются платформы Rentomania, Авито и Арендую.ру. Сравнение функционала данных площадок представлено в таблице \ref{concurent:table}.

\begin{xltabular}{\textwidth}{|X|X|X|X|X|}
\caption{Сравнение существующих решений и аналогов\label{concurent:table}}\\ \hline
Параметр & \centrow Rentomania & \centrow Авито & \centrow Арендую.ру & \centrow Проекти-\ руемая онлайн площадка\\ \hline
\endfirsthead
\continuecaption{Продолжение таблицы \ref{concurent:table}}
~ Параметр & \centrow Rentomania & \centrow Авито & \centrow Арендую.ру & \centrow Проекти-\ руемая онлайн площадка\\ \hline
\finishhead
\leftrow Размещение объявлений арендодателем & \centrow + & \centrow + & \centrow + & \centrow + \\ \hline
\leftrow Поиск и фильтрация товаров & \centrow + & \centrow + & \centrow + & \centrow + \\ \hline
\leftrow Онлайн-бронирование & \centrow + & \centrow -- & \centrow -- & \centrow + \\ \hline
\leftrow Встроенная онлайн-оплата & \centrow + & \centrow + & \centrow -- & \centrow + \\ \hline
\leftrow Система рейтинга и отзывов & \centrow + & \centrow + & \centrow -- & \centrow + \\ \hline
\leftrow Чат между арендатором и арендодателем & \centrow + & \centrow + & \centrow -- & \centrow + \\ \hline
\leftrow Адаптивная вёрстка для мобильных устройств & \centrow + & \centrow + & \centrow -- & \centrow + \\ \hline
\leftrow Верификация пользователей & \centrow + & \centrow -- & \centrow -- & \centrow + \\ \hline
\leftrow Панель аналитики для арендодателя & \centrow -- & \centrow -- & \centrow -- & \centrow + \\ \hline
\leftrow Понятный интерфейс & \centrow -- & \centrow + & \centrow -- & \centrow + \\ \hline
\leftrow Уведомления о статусе аренды & \centrow + & \centrow -- & \centrow -- & \centrow + \\ \hline
\end{xltabular}